\section{Related Work}

\subsection{Classical and Early Linguistic Steganography}

Classical steganography established the distributional perspective that remains central today~\cite{cox2005information}. In linguistic channels, early methods often relied on cover modification and constrained substitution strategies. Later neural generation methods shifted to coverless generation where hidden bits are embedded during token production~\cite{fang2017generating,yang2019rnnstega,dai2019towards}.

\subsection{Arithmetic-Coding Neural Line}

NLS~\cite{ziegler2019neural} is a key reference for modern arithmetic-coding-based linguistic steganography with large neural language models. It links embedding efficiency and distributional closeness, and reports both statistical and human-facing perspectives.

Shen et al.~\cite{shen2020near} further explores practical imperceptibility-oriented adjustments in this family.

Ghostext directly inherits this arithmetic-coding neural line, but emphasizes protocol determinism and fail-closed engineering semantics.

\subsection{Security-Oriented and Provable Directions}

Subsequent work strengthened formal security framing with alternative constructions and proof strategies~\cite{zhang2021provably,kaptchuk2021language,ding2023discop,dewitt2024perfect}. These contributions are important for theory-facing guarantees, but practical deployment still depends on implementation discipline around synchronization and reproducibility.

Ghostext does not claim to supersede these lines. Instead, it targets a complementary space: making protocol behavior explicit, testable, and auditable for real implementation workflows.

\subsection{OD-Stega and Controlled Utility-Security Tradeoff}

OD-Stega~\cite{huang2026odstega} formalizes entropy-maximizing distribution adjustment under divergence constraints, with strong attention to tokenization mismatch and truncation interactions. This is the most relevant direct reference for discussing controlled security relaxation in LLM steganography.

Ghostext uses OD-Stega as a conceptual guide for future extension but currently keeps the base distribution path unchanged. Accordingly, this draft does not claim OD-level utilization improvements.

\subsection{Ghostext's Distinct Contribution in This Landscape}

Relative to prior literature, Ghostext's present contribution is implementation-centric: authenticated opaque packetization, key separation, prompt-and-backend fingerprint binding, deterministic integer quantization, explicit candidate stability filtering, bounded retry/failure controls, and claim-evidence traceability workflow.

These properties are not replacements for detector benchmarks or formal proofs. They are foundations that make later security claims more defensible and reproducible.
